\section{Related Work}
\label{Sec:Related}
%Loop summarization and loop reachability analysis are active areas of research. 
%We discuss techniques related to \solveCHC below.
Termination and nontermination analysis are very active areas of research. We will discuss techniques,
related to Interpolation-based Termination Analysis below.

\textbf{Nontermination.} Nontermination analysis is typically centered around the construction of recurrent sets — sets of states from which execution cannot 
escape~\cite{DBLP:conf/sas/Ben-AmramDG19,DBLP:conf/cade/FrohnG22,DBLP:conf/tacas/BrockschmidtCIK16,DBLP:conf/tacas/ChenCFNO14,DBLP:conf/foveoos/BrockschmidtSOG11,DBLP:conf/cav/LarrazNORR14}.
These techniques can vary from the template-based synthesis to the safety-check based recurrent set construction.
Another traditional approach is safety-based nontermination analysis~\cite{DBLP:journals/entcs/BiereAS02},
that detects nontermination if the same state is reachable within TS more then ones - detecting a lasso.
Newly introduced nontermination technique is the construction of geometric nontermination argument~\cite{DBLP:conf/tacas/LeikeH18} -
a finite representation of an infinite execution that via a sum of geometric series.

\textbf{Termination.} There exist many different techniques for the termination analysis of software.
One of the most researched and applied techniques is the synthesis of the ranking function~\cite{DBLP:conf/vmcai/PodelskiR04,DBLP:conf/cav/KuraUH20,DBLP:journals/jacm/Ben-AmramG14,DBLP:conf/cav/HeizmannHP14,DBLP:journals/jar/GieslABEFFHOPSS17,DBLP:conf/esop/SaritaSGSV26,DBLP:conf/sas/Ben-AmramDG19,DBLP:conf/tacas/UrbanGK16}. 
State of the art techniques are centered on the construction of multiphase or lexicographic ranking functions.
These complex techniques significantly broaden the class of programs for which the termination can be proven.
However, the complexity of transition system can make construction of ranking function difficult.
Other technique, which allows to efficiently prove termination are fixpoint computation~\cite{DBLP:journals/pacmpl/UnnoTGK23}, which
relies on the reduction of termination/nontermination problems into first order fixpoint formula and checking its validity.

There also exists another set of techniques that are based on the construction of well-founded transition
invariants~\cite{TDBLP:conf/lics/PodelskiR04,DBLP:conf/sas/CookPR05,DBLP:conf/cav/KroeningSTW10,DBLP:conf/tacas/TsitovichSWK11}.
These approaches overapproximate the behaviour of the transition system, simplifying the termination proof.
The core problem of these approaches is the generation of transition invariant.
Listed techniques rely either on the template-based generation or on the direct generation of ranking relation based on a particular trace
in the transition system.
In contrast, we propose to do the guided generation of transition invariants, based on the terminating traces and the actual 
termination conditions.
Additionally, our technique is combined with a nontermination approach, increasing the efficiency of both techniques.




